% Part: proof-theory
% Chapter: sequent-calculus
% Section: proof-examples
%READERNOTE{SOL6-A332: संयोग-परिचयाच्या निष्कर्षात दोन पूर्वांग प्रती ठेवून पुढील आकुंचन खरे केले; खोलीच्या आणि स्वयंसिद्धकाच्या दाव्यांचा योग्य scope व निष्फळ संख्यापकाचा अपवाद स्पष्ट केला.}

\documentclass[../../../include/open-logic-section]{subfiles}

\begin{document}

\olfileid{pt}{seq}{ex}

\olsection{सिद्धतांची उदाहरणे}

क्रमवर्तींच्या !!{proof}s तयार करताना बहुतेक
वेळा अंतिम क्रमवर्तीपासून वरच्या दिशेने काम करणे
सोयीचे ठरते: त्यातील एखादे !!a{formula} सक्रिय
!!{formula} म्हणून निवडायचे, त्या क्रमवर्तीच्या वर
संबंधित नियमाचे पूर्वपक्ष लिहायचे, आणि स्वयंसिद्धके
मिळेपर्यंत ही प्रक्रिया पुन्हा करायची. उदाहरणार्थ,
पुढील क्रमवर्ती सिद्ध करायची आहे असे माना:
\[(!C \land !D) \lif !E \Sequent !C \lif (!D \lif !E)\]
$!C \lif (!D \lif !E)$ विचारात घ्या: ते $!A \lif !B$
या रूपाचे असून क्रमवर्तीच्या उजवीकडे आहे. पूर्ण
क्रमवर्तीची \Log{G1c} च्या $\RightR{\lif}$ नियमाशी
जुळवणी करू:
\[\Axiom$ !A, \Gamma \fCenter \Delta, !B$
\RightLabel{\RightR{\lif}}
\UnaryInf$ \Gamma \fCenter \Delta, !A \lif !B $
\DisplayProof\]
यावरून $\Gamma = \{(!C \land !D) \lif !E\}$,
$\Delta =
\emptyset$, $!A \ident !C$ आणि $!B \ident !D \lif !E$
असे घ्यावे लागते. मग सिद्धतेचे शेवटचे निगमन असे असू शकेल:
\[\Axiom$ !C, (!C \land !D) \lif !E \fCenter !D \lif !E$
\RightLabel{\RightR{\lif}}
\UnaryInf$(!C \land !D) \lif !E \fCenter !C \lif (!D \lif !E)$
\DisplayProof\]
हीच कल्पना पुन्हा लावून $!D \lif !E$ ला सक्रिय
सूत्र मानल्यास पुढील मिळते:
\[\Axiom$ !D, !C, (!C \land !D) \lif !E \fCenter !E$
\RightLabel{\RightR{\lif}}
\UnaryInf$!C, (!C \land !D) \lif !E \fCenter !D \lif !E$
\RightLabel{\RightR{\lif}}
\UnaryInf$(!C \land !D) \lif !E \fCenter !C \lif (!D \lif !E)$
\DisplayProof\]
आता $(!C \land !D) \lif !E$ वर लक्ष केंद्रित करून
\Log{G1c} चा $\LeftR{\lif}$ नियम वापरावा लागेल:
\[\Axiom$ \Gamma \fCenter \Delta, !A$
\Axiom$ !B, \Gamma \fCenter \Delta$
\RightLabel{\LeftR{\lif}}
\BinaryInf$ !A \lif !B, \Gamma \fCenter \Delta$
\DisplayProof\]
त्यातून पुढील मिळते:
\[\Axiom$!D, !C \fCenter !E, !C \land !D$
\Axiom$!D, !C, !E \fCenter !E$
\RightLabel{\LeftR{\lif}}
\BinaryInf$ !D, !C, (!C \land !D) \lif !E \fCenter !E$
\RightLabel{\RightR{\lif}}
\UnaryInf$!C, (!C \land !D) \lif !E \fCenter !D \lif !E$
\RightLabel{\RightR{\lif}}
\UnaryInf$(!C \land !D) \lif !E \fCenter !C \lif (!D \lif !E)$
\DisplayProof\]
वरच्या डाव्या क्रमवर्तीतील उजवीकडचे !!{formula}
$!C \land !D$ प्रमुख मानून $\RightR{\land}$ नियम
लावल्यास पुढील मिळते:
\[
\Axiom$!D, !C \fCenter !E, !C$
\Axiom$!D, !C \fCenter !E, !D$
\RightLabel{\RightR{\land}}
\BinaryInf$!D, !C \fCenter !E, !C \land !D$
\Axiom$!D, !C, !E \fCenter !E$
\RightLabel{\LeftR{\lif}}
\BinaryInf$ !D, !C, (!C \land !D) \lif !E \fCenter !E$
\RightLabel{\RightR{\lif}}
\UnaryInf$!C, (!C \land !D) \lif !E \fCenter !D \lif !E$
\RightLabel{\RightR{\lif}}
\UnaryInf$(!C \land !D) \lif !E \fCenter !C \lif (!D \lif !E)$
\DisplayProof
\]
आतापर्यंत सगळे ठीक आहे. आपण वापरलेले नियम
\Log{G1c} आणि~\Log{G3c} मध्ये सारखे असल्याने हे
सर्व \Log{G3c} मध्येही चालते. या !!a{proof} च्या
अपूर्ण भागातील वरच्या सर्व क्रमवर्तींच्या डावीकडे
आणि उजवीकडे तेच !!{formula} येते; पण त्या अद्याप
अगदी स्वयंसिद्धके झालेल्या नाहीत.

\Log{G1c} मध्ये वरच्या क्रमवर्तींच्या
!!{proof}s $!A \Sequent !A$ या रूपाच्या स्वयंसिद्धकांपासून
द्याव्या लागतात. दुर्बलीकरण नियमांनी ते सहज शक्य
आहे. उदाहरणार्थ, सर्वांत डावीकडची वरची क्रमवर्ती
$!D, !C \Sequent
!E, !C$ पुढीलप्रमाणे !!{prove} करता येते:
\[
\Axiom$!C \fCenter !C$
\RightLabel{\LeftR{\Weakening}}
\UnaryInf$!D, !C \fCenter !C$
\RightLabel{\RightR{\Weakening}}
\UnaryInf$!D, !C \fCenter !E, !C$
\]
एकत्रितपणे, \Log{G1c} मधील !!a{proof} अशी दिसते:
\[
\Axiom$!C \fCenter !C$
\RightLabel{\LeftR{\Weakening}}
\UnaryInf$!D, !C \fCenter !C$
\RightLabel{\RightR{\Weakening}}
\UnaryInf$!D, !C \fCenter !E, !C$
\Axiom$!D \fCenter !D$
\RightLabel{\LeftR{\Weakening}}
\UnaryInf$!D, !C \fCenter !D$
\RightLabel{\RightR{\Weakening}}
\UnaryInf$!D, !C \fCenter !E, !D$
\RightLabel{\RightR{\land}}
\BinaryInf$!D, !C \fCenter !E, !C \land !D$
\Axiom$!E \fCenter !E$
\RightLabel{\LeftR{\Weakening}}
\UnaryInf$!C, !E \fCenter !E$
\RightLabel{\LeftR{\Weakening}}
\UnaryInf$!D, !C, !E \fCenter !E$
\RightLabel{\LeftR{\lif}}
\BinaryInf$ !D, !C, (!C \land !D) \lif !E \fCenter !E$
\RightLabel{\RightR{\lif}}
\UnaryInf$!C, (!C \land !D) \lif !E \fCenter !D \lif !E$
\RightLabel{\RightR{\lif}}
\UnaryInf$(!C \land !D) \lif !E \fCenter !C \lif (!D \lif !E)$
\DisplayProof
\]
या सिद्धतेची रचना (विशेषतः तिची उंची) !!{formula}s
$!C$, $!D$, आणि~$!E$ काय आहेत यावर अवलंबून नाही.
वरील !!{proof} मध्ये या चिन्हांच्या जागी कोणतीही
!!{formula}s घातली तरी मिळणारी रचना~\Log{G1c}
मधील !!{proof} राहते. \Log{G1c} मधील ही
सिद्धता \emph{रूपरेषात्मक} आहे.

\Log{G3c} मध्ये स्वयंसिद्धके
$!A, \Gamma \Sequent
\Delta, !A$ या रूपाची क्रमवर्ती असतात; त्यात~$!A$
अण्वीय असले पाहिजे. म्हणून $!D, !C \Sequent
!E, !C$ ही \Log{G3c} ची स्वयंसिद्धकासारखी दिसत
असली ($\Gamma = \{!D\}$; $\Delta =
\{!E\}$), तरी~$!C$ प्रत्यक्षात अण्वीय असेल तेव्हाच
ती \emph{खरोखर} स्वयंसिद्धक असते. $!C$, $!D$, आणि $!E$ ही अण्वीय !!{formula}s
असतील, तर हा सिद्धता-वृक्ष \Log{G3c} मध्ये पूर्ण
सिद्धता असतो. ही पुरेशी अट आहे; इतर बाजूच्या सूत्रामुळे
एखादी वरची क्रमवर्ती स्वयंसिद्धक असण्याची शक्यताही आहे.

काटेकोरपणे पाहता, आपण अजून
\[\Log{G3c} \Proves (!C \land !D) \lif !E \fCenter !C \lif (!D \lif
!E),\] हे दाखवलेले नाही. तरी व्यवहारात एवढेच
करणे आपल्याला पुरेसे आहे (आणि शक्यही आहे).
कारण $!A, \Gamma \Sequent \Delta, !A$ या रूपाची
प्रत्येक क्रमवर्ती \Log{G3c} मध्ये सिद्धतायोग्य
असते, जरी~$!A$ स्वतः अण्वीय नसले तरी.
हे~$!A$ च्या खोली~$\depth{!A}$ वर विगमन करून
दाखवता येते.

\begin{defn}\ollabel{defn:depth}
!!a{formula} ची \emph{खोली}~$\depth{!A}$ पुढीलप्रमाणे
परिभाषित करतो:
\begin{enumerate}
  \item $!A$ अण्वीय असेल तर $\depth{!A} = 0$.
  \item $!A \ident \lnot !B$, $!A \ident \lforall[x][!B]$,
  किंवा $!A \ident \lexists[x][!B]$ असेल तर
  $\depth{!A} = \depth{!B} + 1$.
  \item $!A \ident (!B \land !C)$, $(!B \lor !C)$, किंवा $(!B \lif
  !C)$ असेल तर $\depth{!A} = \max(\depth{!B},\depth{!C}) + 1$.
\end{enumerate}
\end{defn}

कोणत्याही पद~$t$ साठी $\depth{A(x)} =
\depth{A(t)}$ सिद्ध करणे कठीण नाही. आपल्याला
महत्त्वाचे हे की तार्किक नियमात सक्रिय विघटक सूत्रांपेक्षा प्रमुख
!!{formula} ची खोली नेहमी अधिक असते. संरचनात्मक
नियमांसाठी असा दावा नाही; तसेच संख्यापकाच्या काही
\Log{G3c} नियमांत जपलेली प्रमुख आस्थिती स्वतः त्याच
खोलीची असते. उदाहरणार्थ:
\begin{align*}
\depth{P(x)} = \depth{P(a)} & = 0,\\
\depth{\lforall[x][P(x)]} & = 1,\\
\depth{(\lforall[x][P(x)] \lif P(a))} & = \max(1,0) + 1 = 2.
\end{align*}
याचा वापर करून~$\depth{!A}$ वर विगमनाने काही
निष्कर्ष सिद्ध करता येतात; पुढील निष्कर्ष त्यांपैकी एक.

\begin{prop}\ollabel{prop:G1c-axioms}
कोणत्याही~$!A$ साठी, $!A \Sequent !A$ ची अशी
$\Log{G1c}$-!!a{proof} आहे की तिच्यातील सर्व
स्वयंसिद्धके अण्वीय आहेत.
\end{prop}

\begin{proof}
  $\depth{!A}$ वर विगमन करू. $\depth{!A} = 0$
  असेल तर~$!A$ अण्वीय आहे आणि $!A \Sequent !A$
  हीच अपेक्षित रूपाची \Log{G1c} मधील !!{proof} आहे.
  $\depth{!A} > 0$ असेल तर~$!A$ कमी खोलीच्या
  !!{formula}s पासून बनते, उदा., $!A \ident \lnot !B$,
  $!A \ident
  (!B \land !C)$, किंवा $!A \ident \lforall[x][A(x)]$.
  विगमनात्मक गृहीतक असे: $\depth{!D} < \depth{!A}$
  असलेल्या प्रत्येक~$!D$ साठी अण्वीय स्वयंसिद्धकांपासून
  $\Log{G1c} \Proves !D \Sequent !D$.

$!A \ident \lnot !B$ असेल, तर
  $\depth{!B} < \depth{!A}$; विगमनात्मक
  गृहीतकानुसार अण्वीय स्वयंसिद्धकांपासून
  $!B \Sequent !B$ ची \Log{G1c} मधील
  !!a{proof}~$\pi_1$ आहे. तिचा पुढीलप्रमाणे
  $\lnot !B \Sequent \lnot !B$ च्या !!a{proof}
  पर्यंत विस्तार करता येतो:
\[
  \AxiomC{}
  \RightLabel{$\pi_1$}
  \Deduce$!B \fCenter !B$
  \RightLabel{\LeftR{\lnot}}
  \UnaryInf$\lnot !B, !B \fCenter $
  \RightLabel{\RightR{\lnot}}
  \UnaryInf$\lnot !B \fCenter \lnot !B$
  \DisplayProof
\]

$!A \ident (!B \land !C)$ असेल, तर
  $\depth{!B} < \depth{!A}$ आणि
  $\depth{!C} < \depth{!A}$. विगमनात्मक गृहीतकानुसार
  अण्वीय स्वयंसिद्धकांपासून $!B
  \Sequent !B$ आणि $!C \Sequent !C$ यांच्या
  \Log{G1c} मधील !!{proof}s~$\pi_1$ आणि~$\pi_2$
  आहेत. त्यांचा पुढीलप्रमाणे
  $!B \land !C \Sequent !B \land !C$ च्या
  !!a{proof} पर्यंत विस्तार करता येतो:
\[
  \AxiomC{}
  \RightLabel{$\pi_1$}
  \Deduce$!B \fCenter !B$
  \RightLabel{\LeftR{\Weakening}}
  \UnaryInf$!B, !C \fCenter !B$
  \RightLabel{\LeftR{\land}}
  \UnaryInf$!B \land !C, !C\fCenter !B$
  \RightLabel{\LeftR{\land}}
  \UnaryInf$!B \land !C, !B \land !C \fCenter !B$
  \AxiomC{}
  \RightLabel{$\pi_2$}
  \Deduce$!C \fCenter !C$
  \RightLabel{\LeftR{\Weakening}}
  \UnaryInf$!B, !C \fCenter !C$
  \RightLabel{\LeftR{\land}}
  \UnaryInf$!B \land !C, !C \fCenter !C$
  \RightLabel{\LeftR{\land}}
  \UnaryInf$!B \land !C, !B \land !C \fCenter !C$
  \RightLabel{\RightR{\land}}
  \BinaryInf$!B \land !C, !B \land !C \fCenter !B \land !C$
  \RightLabel{\LeftR{\Contraction}}
  \UnaryInf$!B \land !C \fCenter !B \land !C$
  \DisplayProof
  \]

आता $!A \ident \lforall[x][B(x)]$ असे माना.
  $\lforall[x][B(x)]$ मध्ये न येणारा
  !!a{constant}~$c$ घ्या. मग $\depth{B(c)} =
  \depth{!B(x)} < \depth{!A}$; विगमनात्मक गृहीतकानुसार
  अण्वीय स्वयंसिद्धकांपासून $!B(c)
  \Sequent !B(c)$ ची \Log{G1c} मधील
  !!a{proof}~$\pi_1$ आहे. तिचा पुढीलप्रमाणे
  $\lforall[x][B(x)] \Sequent \lforall[x][B(x)]$
  च्या !!a{proof} पर्यंत विस्तार करता येतो:
\[
  \AxiomC{}
  \RightLabel{$\pi_1$}
  \Deduce$!B(c) \fCenter !B(c)$
  \RightLabel{\LeftR{\lforall}}
  \UnaryInf$\lforall[x][!B(x)] \fCenter !B(c)$
  \RightLabel{\RightR{\lforall}}
  \UnaryInf$\lforall[x][!B(x)] \fCenter \lforall[x][!B(x)]$
  \DisplayProof
\]

उरलेले प्रसंग प्रश्न म्हणून सोडले आहेत.
\end{proof}

\begin{prob}
  $!A \ident (!B \lor !C)$, $!A \ident (!B \lif !C)$
  आणि $!A \ident \lexists[x][B(x)]$ या प्रसंगांचे
  काम करून \olref{prop:G1c-axioms} ची सिद्धता पूर्ण करा.
\end{prob}

आपल्या कामासाठी $!A \ident \lnot !B$ या
प्रसंगात पुढील !!{proof} ही वापरता आली असती:
\[
\AxiomC{}
\RightLabel{$\pi_1$}
\Deduce$!B \fCenter !B$
\RightLabel{\RightR{\lnot}}
\UnaryInf$\fCenter !B, \lnot !B$
\RightLabel{\LeftR{\lnot}}
\UnaryInf$\lnot !B \fCenter \lnot !B$
\DisplayProof
\]
मात्र काही क्रमवर्ती पद्धतींमध्ये हे चालले नसते.
अंतःप्रज्ञावादी क्रमवर्ती कलन~\Log{G1i} हे
\Log{G1c} सारखेच आहे; फरक इतकाच की प्रत्येक
क्रमवर्तीच्या उजव्या बाजूस जास्तीत जास्त एकच
!!{formula} ठेवता येते. $\Delta = \emptyset$
असेल, तर पहिली !!{proof} ही~\Log{G1i} मध्येही
!!a{proof} असेल; परंतु दुसरी नसेल, कारण तिच्या
एका क्रमवर्तीच्या उजवीकडे $!B$ आणि~$\lnot !B$
दोन्ही आहेत.

$!A \ident \lforall[x][!B(x)]$ असताना
\Log{G1c} मध्येही नियमांचा दुसरा क्रम वापरता
आला नसता. पुढील रचना \emph{!!a{proof} नाही}:
\[
  \AxiomC{}
  \RightLabel{$\pi_1$}
  \Deduce$!B(c) \fCenter !B(c)$
  \RightLabel{\RightR{\lforall}}
  \UnaryInf$!B(c) \fCenter \lforall[x][!B(x)]$
  \RightLabel{\LeftR{\lforall}}
  \UnaryInf$\lforall[x][!B(x)] \fCenter \lforall[x][!B(x)]$
  \DisplayProof
\]
कारण त्यात~$\RightR{\lforall}$ ची आयगेन चराची
अट मोडली जाते.
येथे~$x$ हे~$!B(x)$ मध्ये खरोखर मुक्त आहे असे मानतो,
म्हणून~$c$ डावीकडील~$!B(c)$ मध्ये येतो. संख्यापक
निष्फळ असेल आणि~$c$ त्या सूत्रात येतच नसेल, तर
दाखवलेल्या आयगेन-चर अटीचे उल्लंघन होत नाही.

\begin{prop}\ollabel{prop:G3c-axioms}
$!A, \Gamma
\Sequent \Delta, !A$ या रूपाची कोणतीही क्रमवर्ती
(~$!A$ \emph{अण्वीय} असण्याची आवश्यकता नाही)
\Log{G3c} मध्ये सिद्धतायोग्य आहे.
\end{prop}

\begin{proof}
ही सिद्धता \olref{prop:G1c-axioms} च्या सिद्धतेसारखीच
आहे, पण विगमनात्मक गृहीतकाबाबत काळजी घ्यावी
लागते. $n =
\depth{!A} = 0$ हा प्रसंग पुन्हा सोपा आहे.
$n=0$ असेल, तर~$!A$ अण्वीय आहे, आणि $!A,
\Gamma \Sequent \Delta, !A$ हे~\Log{G3c}
चे स्वयंसिद्धक आहे.

आता $n>0$ असे माना. पुन्हा प्रसंग वेगळे करू.
विगमनात्मक गृहीतक असे: $\depth{!D} < n$ असलेल्या
सर्व~$!D$ साठी आणि सर्व~$\Pi$ व~$\Lambda$
साठी $\Log{G3c} \Proves !D, \Pi \Sequent \Lambda, !D$.
$!A \ident (!B \land !C)$ हा प्रसंग पाहू. विगमनात्मक
गृहीतक दोनदा वापरू: एकदा $!D = !B$,
$\Pi = !C, \Gamma$, आणि $\Lambda
= \Delta$ घेऊन $!B, !C, \Gamma \Sequent
\Delta, !B$ ची !!a{proof}~$\pi_1$ मिळते; आणि
दुसऱ्यांदा $!D = !C$, $\Pi = !B, \Gamma$, आणि
$\Lambda = \Delta$ घेऊन $!B, !C, \Gamma \Sequent
\Delta, !C$ ची !!a{proof}~$\pi_2$ मिळते.
$!B \land !C, \Gamma \Sequent
\Delta, !B \land !C$ ची !!{proof} पुढीलप्रमाणे:
\[
\AxiomC{}
\RightLabel{$\pi_1$}
\Deduce$!B, !C, \Gamma \fCenter \Delta, !B$
\RightLabel{\LeftR{\land}}
\UnaryInf$!B \land !C, \Gamma \fCenter \Delta, !B$
\AxiomC{}
\RightLabel{$\pi_2$}
\Deduce$!B, !C, \Gamma \fCenter \Delta, !C$
\RightLabel{\LeftR{\land}}
\UnaryInf$!B \land !C, \Gamma \fCenter \Delta, !C$
\RightLabel{\RightR{\land}}
\BinaryInf$!B \land !C, \Gamma \fCenter \Delta, !B \land !C$
\DisplayProof
\]

$!A \ident \lforall[x][!B(x)]$ साठी
~$\lforall[x][!B(x)]$, $\Gamma$ किंवा~$\Delta$
मध्ये न येणारा !!a{constant} निवडा.
$!D = !B(c)$, $\Pi = \lforall[x][!B(x)], \Gamma$
आणि $\Lambda = \Delta$ घेऊन विगमनात्मक गृहीतक
लावा. त्यातून $!B(c),
\lforall[x][!B(x)], \Gamma \Sequent \Delta, !B(c)$
ची !!a{proof}~$\pi_1$ मिळते; तिच्यापासून:
\[
\AxiomC{}
\RightLabel{$\pi_1$}
\Deduce$!B(c), \lforall[x][!B(x)], \Gamma \fCenter \Delta, !B(c)$
\RightLabel{\LeftR{\lforall}}
\UnaryInf$\lforall[x][!B(x)], \Gamma \fCenter \Delta, !B(c)$
\RightLabel{\RightR{\lforall}}
\UnaryInf$\lforall[x][!B(x)], \Gamma \fCenter \Delta, \lforall[x][!B(x)]$
\DisplayProof
\]
आपण~$c$ ची केलेली निवड~$\RightR{\lforall}$
ची आयगेन चराची अट पूर्ण करते.
\end{proof}

\begin{prob}
  $!A \ident (!B \lif !C)$ आणि $!A \ident
  \lexists[x][B(x)]$ या प्रसंगांचे काम करून
  \olref{prop:G3c-axioms} ची सिद्धता पुढे न्या.
\end{prob}

\end{document}
